\begin{table}[p]
\centering
\small
\setlength{\tabcolsep}{4pt}
\caption{Reference-wise relative improvement (\%) of GC-\(0.125\) over D-MMM for the low-density-region (LDR) Hellinger score at \(\ell_{\mathrm{LDR}}=0.01\) in the perfect-model experiment. Percentages are computed separately for each held-out reference and target period from area-weighted LDR Hellinger scores, as in Table~\ref{tab:pme_reference_improvement}. Positive values favour GC-\(0.125\). Summary rows give the mean, median, and interquartile range (Q1--Q3) across 15 equally weighted reference models.}
\label{tab:pme_reference_improvement_ldr}
\begin{tabular}{lrrrrr}
\toprule
Held-out reference & 1976--2000 & 2001--2025 & 2026--2050 & 2051--2075 & 2076--2100 \\
\midrule
ACCESS-CM2 & 41.5 & 41.7 & 40.4 & 38.2 & 36.5 \\
BCC-CSM2-MR & 45.4 & 46.0 & 45.3 & 42.7 & 29.0 \\
CMCC-CM2-SR5 & 49.0 & 49.5 & 50.0 & 48.9 & 47.5 \\
CNRM-CM6-1-HR & 43.0 & 42.0 & 45.0 & 41.6 & 29.8 \\
CanESM5 & 43.0 & 41.8 & 38.7 & 33.7 & 27.7 \\
EC-Earth3-CC & 40.1 & 41.6 & 40.8 & 38.7 & 35.9 \\
GFDL-ESM4 & 36.8 & 35.9 & 34.1 & 33.2 & 28.5 \\
INM-CM5-0 & 40.5 & 39.3 & 36.6 & 31.9 & 22.8 \\
MIROC-ES2L & 58.0 & 58.1 & 57.9 & 56.8 & 56.8 \\
MIROC6 & 58.2 & 59.5 & 59.3 & 60.6 & 59.5 \\
MPI-ESM1-2-LR & 41.2 & 40.6 & 37.1 & 36.3 & 35.0 \\
MRI-ESM2-0 & 47.5 & 46.7 & 46.4 & 40.7 & 32.8 \\
NESM3 & 40.9 & 39.1 & 36.5 & 34.6 & 33.9 \\
NorESM2-MM & 43.0 & 41.7 & 38.0 & 33.0 & 29.8 \\
TaiESM1 & 42.8 & 44.4 & 47.3 & 46.6 & 45.3 \\
\midrule
Mean & 44.7 & 44.5 & 43.5 & 41.2 & 36.7 \\
Median & 43.0 & 41.8 & 40.8 & 38.7 & 33.9 \\
IQR (Q1--Q3) & 41.1--46.5 & 41.1--46.4 & 37.5--46.8 & 34.1--44.7 & 29.4--40.9 \\
\bottomrule
\end{tabular}
\end{table}
